\section[A new formula for the Weyl group generators on any integrable module]{A new formula for the Weyl group generators\\ on any integrable module}

It is remarkable that the proof of the following theorem, starting from a simple calculation in the two-dimensional representation of
$\fsl(2,\C)_i$, generalizes to a result in any integrable module for any Kac--Moody algebra $\fg$. Special cases of this formula appeared
in some physics papers, for example, in \cite{DamourHillmann09}, where Damour and Hillmann found it in a representation of
$K(E_{10})$.

It is known \cite{Kac90} that for any integrable representation $\phi\colon \fg\to {\rm End}(V)$, the group $\tW$ generated by the operators
\[\tw_i^\phi = \exp(\phi(e_i)) \exp(\phi(-f_i) \exp(\phi(e_i)),\qquad 1\leq i\leq n\]
is a subgroup of the Kac--Moody group $G$ that acts on $V$ as follows. If $V_\mu$ is the $\mu$-weight space of $V$, then $\tw_i^\phi (V_\mu) = V_{w_i(\mu)}$, where $w_i$ is the reflection with respect to the simple root $\alpha_i$ in the dual space, $\fh^*$, of the standard Cartan subalgebra, $\fh$. The Weyl group, $W$, is the Coxeter group generated by those simple reflections, and is also defined on $\fh$ by the formula $w_i(h) = h - \alpha_i(h) h_i$. There is a surjection from $\tW$ onto $W$ whose kernel is given by Remark 3.8 in \cite{Kac90}. In the case when $\phi$ is the adjoint representation, the restriction of each $\tw_i^\phi$ to $\fh$ equals $w_i$.

For any integrable representation $\phi\colon \fg\to {\rm End}(V)$, define the ``compact" operators from $G^\theta$,
\[s_i^\phi = \exp(\phi(\pi x_i)) = \exp(\phi(\pi(e_i - f_i)/2)),\qquad 1\leq i\leq n.\]
We prove that $\tw_i^\phi = s_i^\phi$. This provides a new description of $\tW$ as generated by operators from the real compact form. This generalizes a result discovered from a physics point of view by Damour and Hillmann in a representation of~$K(E_{10})$.

\begin{Theorem}\label{theorem:newWeylFormula} Let $\fg$ be any Kac--Moody algebra of rank $n$ with the usual Chevalley generators $e_i$, $f_i$, $h_i$ and let $\phi\colon \fg\to {\rm End}(V)$ be any integrable representation. Then for $1\leq i\leq n$, we have
\[\exp(\phi(e_i)) \exp(\phi(-f_i) \exp(\phi(e_i)) = \exp(\phi(\pi(e_i - f_i)/2)).\]
\end{Theorem}

\begin{proof}Using the notations as above, for $1\leq i\leq n$, let $\fg_i = \fsl(2,\C)_i$ be the Lie subalgebra of~$\fg$ with basis~$\{e_i,f_i,h_i\}$. Then $V$ has a direct sum decomposition into irreducible $\fg_i$-modules
\[V = \bigoplus_{j\in J} V_j(m),\]
where $\dim(V_j(m)) = m+1$ and the index set $J$ includes information about the $\fg_i$ highest weight vector in $V_j(m)$ that locates it in~$V$.
It is clear that $\tw_i^\phi = s_i^\phi$ on $V$ if and only if their restrictions to each $V_j(m)$ are equal. They are certainly equal for all trivial one-dimensional modules, $V_j(0)$. On any two-dimensional module, $V_j(1)$, the computation comes down to the simple fact that
\[\exp\left( \bm 0&1\\0&0\ebm \right) \exp\left( \bm 0&0\\-1&0\ebm \right) \exp\left( \bm 0&1\\0&0\ebm \right)
= \bm 1&1\\0&1\ebm \bm 1&0\\-1&1\ebm \bm 1&1\\0&1\ebm = \bm 0&1\\-1&0\ebm \]
matches
\[\exp\left( \bm 0&t\\-t&0\ebm \right) = \bm \cos(t)&\sin(t) \\-\sin(t)&\cos(t) \ebm,\]
when $t = \pi/2$. The well-known tensor product decomposition of irreducible $\fsl(2,\C)$-modules says that for any integer $m\geq 1$, we have
\[V(m)\otimes V(1) = V(m+1) \oplus V(m-1).\]
So if the operators are equal on modules $V(m)$ and $V(1)$ then they are equal on the tensor product, so they are equal on the component
$V(m+1)$. This proves by induction that they are equal on any $V_j(m)$ that occurs in $V$.
\end{proof}
